% =============================================================================
% From Eclipse Demon to Lunar Apogee
% Companion paper — pdflatex + CJKutf8, two-column
% =============================================================================
\documentclass[10pt,a4paper]{article}

% --- Packages (matches when-the-stars-disagree.tex) -------------------------
\usepackage[utf8]{inputenc}
\usepackage[T1]{fontenc}
\usepackage{CJKutf8}
\usepackage{amsmath,amssymb,amsfonts}
\usepackage{booktabs}
\usepackage{multirow}
\usepackage{graphicx}
\usepackage{xcolor}
\usepackage{hyperref}
\usepackage[margin=0.75in]{geometry}
\usepackage{enumitem}
\usepackage{float}
\usepackage{caption}
\usepackage{subcaption}
\usepackage{siunitx}
\usepackage{longtable}
\usepackage[numbers,sort&compress]{natbib}
\usepackage{xurl}
\usepackage{authblk}
\usepackage{pgfplots}
\usepackage{tikz}
\usepackage{gensymb}
\usepackage{adjustbox}
\usepackage[framemethod=tikz]{mdframed}

\pgfplotsset{compat=1.18}

\setlength{\columnsep}{18pt}
\setlength{\emergencystretch}{3em}

\hypersetup{
    colorlinks=true,
    linkcolor=blue!70!black,
    citecolor=green!50!black,
    urlcolor=blue!60!black,
}

% CJK helper
\newcommand{\zh}[1]{\begin{CJK}{UTF8}{bsmi}#1\end{CJK}}

% ORCID helper
\newcommand{\orcidlink}[1]{%
  \href{https://orcid.org/#1}{\textcolor{green!50!black}{\textsf{iD}}}%
}

% pgfplots colors
\definecolor{apogeeblue}{HTML}{2166AC}
\definecolor{nodered}{HTML}{D6604D}

% --- Title -------------------------------------------------------------------

\title{%
  \textbf{From Eclipse Demon to Lunar Apogee:}\\[4pt]
  \large Six Modes of Transformation in Cross-Civilizational\\
  Astronomical Transmission%
}

\author[1]{Albert Kwun Tai Hui\,\orcidlink{0009-0000-8978-0231}}
\affil[1]{Security Ronin}

\date{March 2026}

% =============================================================================
\begin{document}

\twocolumn[
\maketitle

\begin{abstract}
When astronomical concepts travel across civilizational boundaries, they do not
arrive unchanged. This paper argues that such transformations are systematic
rather than random, and proposes a six-mode theoretical framework---transliteration,
reidentification, indigenous restructuring, hybridization, foreign correction, and
algorithmic fossilization---derived from an extended case study spanning thirteen
centuries.

The case study follows the Indian shadow planet Ketu (\zh{計都}) and its coordinate
reference frame as both were absorbed into Chinese \zh{七政四餘} (Seven Governors,
Four Remainders) astrology.  On the coordinate axis the paper traces four
successive regime changes: ecliptic sidereal coordinates in the Song dynasty
(\zh{三辰通載}); equatorial sidereal in the Yuan (\zh{鄭氏星案}); a novel
quasi-ecliptic sidereal hybrid in the Ming (\zh{果老星宗}); and ecliptic tropical
under the Jesuit reform of 1645.  In parallel, Ketu's astronomical identity was
quietly reidentified---from the descending lunar node to the osculating lunar
apogee---a shift demonstrated by Niu Weixing's 1994 recomputation of Tang-dynasty
ephemeris data \cite{niu1995ketu}, before being reverted by Schall von Bell in the
\textit{Shixian} calendar.

The paper contributes: (1)~computational verification of the apsidal identification
using modern JPL ephemeris algorithms; (2)~the first rigorous definition of the
\zh{似黃道} (quasi-ecliptic) coordinate system; and (3)~the six-mode framework as
a reusable instrument for the comparative history of transmitted astronomy.
\end{abstract}
\vspace{12pt}
]

% === §1 Introduction =========================================================
\section{Introduction}\label{sec:intro}
% ~1,500 words — Task 6

\subsection*{The Silk Road as a Channel for Astronomical Ideas}

Astronomical knowledge traveled the Silk Road alongside silk, spices, and
religious texts.  Between the second and ninth centuries~CE, Greek planetary
models, Indian \textit{siddh\=anta} tables, Persian \textit{z\={\i}j} manuals,
and Babylonian omen compendia all made their way into Central Asia and China via
a network of court astronomers, Buddhist missionaries, Sogdian merchants, and
imperial tribute missions.  Unlike a bolt of silk, however, an astronomical concept
does not arrive at its destination intact.  It is translated into a new language,
fitted into an existing institutional structure, combined with native techniques,
challenged by fresh observations, and eventually enshrined in software whose users
have forgotten what it ever meant.  These transformations leave detectable traces
in historical texts---and in the computed positions of celestial bodies.

This paper argues that the transformations are systematic rather than random.
Using a single case study that spans thirteen centuries and four Chinese dynasties,
we identify six recurring patterns---six \emph{modes} of transformation---that
structured the reception of Indian astronomy in China.  The case study concerns
two related objects: the Indian shadow planet Ketu (\zh{計都}) and the coordinate
reference frame used to specify its position.  Both entered Chinese astronomy in
718~CE through the translation of the \textit{Jiuzhi li} (\zh{九執曆},
\textit{Nine Luminaries Calendar}) by Gautama Siddh\=artha (\zh{瞿曇悉達}), a
scholar of Indian descent serving at the Tang court.  Both underwent a remarkable
sequence of transformations over the following millennium, and both continue to
influence popular Chinese-language astrology software today.

\subsection*{Two Parallel Threads}

The paper traces two intertwined threads that run from 718~CE to the present.

\textit{Thread~1: The coordinate reference frame.}  Every astronomical position
requires a reference frame---a choice of origin, fundamental plane, and direction
of increasing longitude.  The original Indian system used ecliptic coordinates
measured from a sidereal zero point (a fixed star rather than the moving
vernal equinox).  Chinese astronomy in the Tang and Song dynasties preferred
equatorial coordinates referenced to the twenty-eight lunar lodges (\zh{二十八宿},
\textit{xiu}).  When Ketu entered the Chinese tradition, a tension arose between
the imported ecliptic frame and the native equatorial one.  That tension was
resolved---and re-resolved---four times:
\begin{enumerate}[leftmargin=*, label=(\arabic*)]
  \item \textbf{Song ecliptic sidereal.}  The \zh{三辰通載} (c.~11th century~CE)
    presents Ketu positions in a coordinate frame that is recognisably ecliptic
    and sidereal, close to the Indian original.
  \item \textbf{Yuan equatorial sidereal.}  The \zh{鄭氏星案} (late Yuan, c.~14th
    century) converts all positions into equatorial coordinates aligned with the
    native lodge system---a thoroughgoing indigenization of the imported framework.
  \item \textbf{Ming quasi-ecliptic sidereal.}  The \zh{果老星宗} (Ming dynasty,
    c.~15th--16th century) introduces what we term the \zh{似黃道} (quasi-ecliptic)
    frame: a hybrid that applies ecliptic-style longitude increments along a plane
    that is neither the true ecliptic nor the equator.  This is the least-studied
    coordinate system in Chinese historical astronomy, and the paper offers the
    first formal definition.
  \item \textbf{Qing ecliptic tropical.}  The Jesuit-led calendar reform of 1645
    (\textit{Shixian li}, \zh{時憲曆}) reimposed European tropical ecliptic
    coordinates, resetting the accumulated drift and discarding the Ming hybrid.
\end{enumerate}
The trajectory is not a smooth evolution but a sequence of reversals---``there
and back again, but not quite'': the Qing frame resembles the Song one in being
ecliptic, but has shifted from sidereal to tropical reckoning.

\textit{Thread~2: The identity of Ketu.}  In Indian astronomy Ketu (\textit{ketu},
literally ``comet'' or ``banner'') designates the descending lunar node---one of
the two points where the Moon's orbit crosses the ecliptic.  When the Moon crosses
the descending node moving southward, conditions favour a lunar eclipse; when
it crosses the ascending node Rahu (\zh{羅睺}), conditions favour a solar eclipse.
The two nodes are therefore deeply embedded in eclipse prediction.

Niu Weixing's landmark 1994 analysis of the \zh{七曜攘災訣} (806~CE), a Tang
ephemeris attributed to the monk Amoghavajra's circle, demonstrated that the
Ketu positions recorded therein do \emph{not} track the descending node
\cite{niu1995ketu}.  They track the osculating lunar apogee---the point in the
Moon's elliptical orbit farthest from Earth.  The apogee was known to Indian
astronomers as \textit{mand\=occa} (``slow apex'') and in Persian-influenced
sources as \textit{K\=etus}; its conflation with the node appears to have
occurred during or shortly after the Tang transmission.  By the Ming period,
Ketu as apogee was thoroughly assimilated into \zh{七政四餘} practice, where it
joined Rahu (\zh{羅睺}), the Purple Star (\zh{紫氣}), and the Monthly Comet
(\zh{月孛}) as the four ``remainders''---invisible bodies that governed fate
alongside the seven visible planets.  The Jesuit reform then reverted Ketu to
the descending node, erasing five centuries of accumulated reidentification.

\subsection*{Six Modes of Transformation}

The six modes proposed here are not a priori categories imposed on the data.
They emerge from the case study as the recurring patterns by which astronomical
concepts were altered during transmission.  We state them here in condensed form;
each receives fuller treatment in the sections that follow.

\textbf{Mode~1: Transliteration.}  Foreign technical terms are adopted into the
receiving tradition with minimal semantic change, using phonetic approximation.
\textit{R\=ahu} becomes \zh{羅睺} (\textit{lu\'ohou}); \textit{Ketu} becomes
\zh{計都} (\textit{j\`{\i}d\=u}).  The sounds are borrowed; the astronomical
meanings---initially---travel with them.  Transliteration is the entry mechanism:
it creates a named slot in the receiving system before that slot has been filled
with indigenous content.

\textbf{Mode~2: Reidentification.}  The name persists while the astronomical
referent shifts.  The transformation is invisible at the lexical level but
detectable by comparing computed positions against historical ephemeris records.
Ketu's drift from descending node to osculating apogee is the central instance;
we suspect the shift happened in stages between 718 and 806~CE as Tang
practitioners noticed that the imported node-based algorithms gave increasingly
poor results and sought a better-fitting celestial body.

\textbf{Mode~3: Indigenous restructuring.}  The native tradition actively
overwrites the imported framework by substituting its own coordinate system,
calendar reckoning, or theoretical structure.  The Yuan equatorial conversion
is the clearest example: Yuan astronomers did not argue with the Indian ecliptic
frame; they simply replaced it with the equatorial lodge-coordinate system that
their \textit{Shoushi li} (\zh{授時曆}) reform had just made canonical.

\textbf{Mode~4: Hybridization.}  Elements of the imported and native systems are
combined into a novel structure that belongs fully to neither parent tradition.
The Ming \zh{似黃道} system is the prime instance.  It preserves ecliptic-style
longitude from the Indian source while anchoring the zero-point and plane to a
native star-count system; the result is not Indian ecliptic astronomy and not
Chinese lodge astronomy, but a tertium quid that persists in Chinese-language
software to this day.

\textbf{Mode~5: Foreign correction.}  An external authority---a new wave of
transmission, a conquering dynasty, or an imported expertise---imposes a reset
that overrides accumulated indigenous development.  Schall von Bell's 1645
\textit{Shixian li} reform is the paradigm case.  Backed by Qing imperial
authority, it discarded the Ming quasi-ecliptic frame, reverted Ketu to the
descending node, and replaced sidereal with tropical reckoning---all in a single
edict.

\textbf{Mode~6: Algorithmic fossilization.}  A living astronomical practice,
once encoded in tables or (later) software, becomes separated from its theoretical
foundation.  Users apply the algorithm without access to---or interest in---the
historical reasoning behind it.  Errors accumulate unnoticed because the algorithm
is self-consistent even if its foundations have shifted.  This mode is not
historically confined to the past: modern Chinese-language astrology software
routinely inherits the Ming \zh{似黃道} frame without labelling it as such,
producing systematic positional errors of several degrees when compared against
modern ephemerides.

\subsection*{Method}

The paper combines close reading of primary astronomical texts with computational
verification using modern ephemeris algorithms.  For the computational component
we use the Swiss Ephemeris library (version~2.10), which implements the JPL
DE441 planetary and lunar solution.  This allows us to ask a precise empirical
question: does a given historical Ketu ephemeris track the lunar apogee, the
descending node, or something else?  We convert historical sexagesimal positions
to modern decimal degrees, apply the appropriate calendar and epoch corrections,
and compare against Swiss Ephemeris output for the descending node and for the
osculating apogee at the same Universal Time.  The method follows and extends
the approach used by Niu Weixing \cite{niu1995ketu} for the \zh{七曜攘災訣} and
applied to Ming texts by Niu in his later synthesis \cite{niu2010fourremains}.
Computational details and residual tables are presented in \S\ref{sec:computation}.

\subsection*{Relationship to Existing Scholarship}

This paper builds on four converging research programmes.

Niu Weixing's identification of Ketu as the lunar apogee
\cite{niu1995ketu,niu2010fourremains} is the single most important prior result
for our argument.  Without Niu's recomputation, the reidentification mode would
remain invisible.  Our contribution is to situate Niu's finding within a broader
transformative sequence and to provide independent computational confirmation
using a modern ephemeris unavailable to Niu in 1994.

Bill Mak's work on pseudo-planets traces the historical trajectory of Rahu and
Ketu from Hellenistic and Indian origins into China and Japan, with particular
attention to Persian-mediated transmission \cite{mak2014pseudoplanets,mak2024persian}.
Mak documents the textual genealogy that brought shadow-planet concepts to East
Asia; the present paper takes up where Mak's textual history ends, asking what
happened to those concepts inside China over the long Ming--Qing transition.

Jeffrey Kotyk has produced detailed studies of the sinicization of Indo-Iranian
astrology in the Tang and Song courts \cite{kotyk2018sinicization,kotyk2017buddhist}.
Kotyk establishes the Buddhist monastic networks---particularly Amoghavajra's
school---through which Indian astronomical texts circulated in the eighth century.
Our \S\ref{sec:tang} on the Tang transmission depends heavily on Kotyk's
reconstruction of institutional context.

Liu Ts'un-yan's work on the Jesuit encounter with Chinese astrology
\cite{liu2020missionary} provides the documentary foundation for our account of
Mode~5 (foreign correction) in \S\ref{sec:jesuit}.  The Jesuit reform was not
merely a technical event; it was an act of epistemic authority, and Liu's analysis
of Schall von Bell's rhetorical strategies illuminates why the reversion of Ketu
to the descending node was politically possible in 1645.

A companion paper, ``When the Stars Disagree: A Cross-Civilizational Analysis
of Computational Errors in Astronomical Divination,'' presents a broader taxonomy
of error types arising from coordinate-frame mismatches across the Chinese,
Hellenistic, Indian, and Persian traditions.  The present paper concentrates on
the transformation \emph{mechanisms} rather than the error \emph{outcomes}.

\subsection*{Structure of the Paper}

Section~\ref{sec:indian} describes the Indian original: the conceptual status of
Rahu and Ketu in Sanskrit \textit{siddh\=anta} astronomy and the ecliptic sidereal
coordinate frame in which they were embedded.  Sections~\ref{sec:tang}--\ref{sec:jesuit}
trace the four dynasties and four coordinate regimes in sequence.
Section~\ref{sec:fossil} analyses the ongoing algorithmic fossilization in modern
software.  Section~\ref{sec:computation} presents the computational verification.
Section~\ref{sec:discussion} draws the six modes together into the theoretical
framework and discusses their applicability beyond the Ketu case.
Section~\ref{sec:conclusion} concludes.

% === §2 The Indian Original ===================================================
\section{The Indian Original: Rahu, Ketu, and Ecliptic Sidereal Coordinates}
\label{sec:indian}
% ~1,000 words — Task 7

\subsection*{Ecliptic Sidereal Coordinates}

Indian \textit{Jyotish} astronomy rests on a coordinate framework that differs
from both the Greek tropical system and the native Chinese equatorial tradition.
The fundamental reference plane is the \emph{ecliptic}---the apparent annual path
of the Sun against the background of fixed stars.  Along this plane the sky is
divided into twelve \textit{r\=asi} (zodiac signs), each occupying exactly 30° of
arc, yielding the familiar circle of 360°.  What distinguishes the Indian system
from its Greek counterpart is the anchor point: where Hellenistic and later
Western astronomy measures longitude from the vernal equinox (a moving point
that drifts westward through the stars at roughly 50 arcseconds per year due to
the precession of Earth's rotation axis), Indian astronomy measures from fixed
stars.  A position is therefore \emph{sidereal}---\textit{nirayana}, ``without
the \textit{ayana}'' or solstitial shift---rather than tropical
(\textit{s\=ayana}).

The distinction is consequential.  Over the roughly two millennia since the
Indian and Hellenistic systems last agreed on the position of the vernal equinox
(approximately 285~CE by the most widely used \textit{Lahiri} \textit{ayanamsha},
though estimates range from 285 to 560~CE depending on the \textit{ayanamsha}),
precession has opened a gap of nearly 24° between tropical and sidereal
longitudes.  A planet at 10° tropical Aries sits near 16° sidereal Pisces.
Every translation between Indian and Chinese sources must contend with this
offset \cite{pingree1981jyotih}.

The twelve \textit{r\=asi}---\textit{Mesha} (Aries) through \textit{M\={\i}na}
(Pisces)---serve as the containers within which planetary positions are reported.
They are equal divisions, not constellations of variable extent, so the
\textit{r\=asi} framework is a clean mathematical abstraction.  Each
\textit{r\=asi} is further subdivided into \textit{nakshatra} (lunar mansion)
segments, linking the zodiacal and lunar-mansion systems in a unified
architecture.  It was this architecture---ecliptic plane, 360° circle,
equal-division signs, nine celestial bodies---that Tang dynasty translators would
import wholesale into China in 718~CE.

\subsection*{Rahu and Ketu: The Eclipse Demons}

Within the nine \textit{graha} (``seizers'' or celestial influences) of Indian
astronomy, two occupy a peculiar metaphysical status.  Rahu and Ketu are not
visible planets but mathematical points: the two nodes where the Moon's orbital
plane intersects the ecliptic.  At these intersections, and only at these
intersections, solar and lunar eclipses are geometrically possible.  Because
eclipses were among the most feared celestial events in ancient societies, the
nodes acquired a mythology commensurate with their power.

The origin story is preserved in the \textit{Mah\=abh\=arata} and the
\textit{Bh\=agavata Pur\=ana}.  During the churning of the cosmic ocean the gods
produced \textit{am\d{r}ta}, the nectar of immortality.  A demon named
Svarbh\=anu disguised himself among the gods and drank from the vessel.  The Sun
and Moon recognized the impostor and alerted Vishnu, who immediately severed
Svarbh\=anu's head from his body with his discus.  Because the demon had already
swallowed the nectar, however, both head and tail survived as immortal fragments:
the head became Rahu, the tail became Ketu.  Eternally enraged, they periodically
swallow the Sun and Moon---producing eclipses---only for the luminaries to emerge
from the severed neck moments later.

This mythological account maps precisely onto the astronomical reality.  Rahu,
the ascending (north) lunar node, is where the Moon crosses from south to north
of the ecliptic; Ketu, the descending (south) node, is where it crosses from
north to south.  The two nodes are always exactly 180° apart---head and tail of
the same orbital intersection---and they precess westward around the ecliptic
with a period of approximately 18.6 years.  The identification is astronomically
unambiguous and remained stable across Indian tradition for more than a millennium
\cite{pingree1981jyotih}.

One complication deserves note before this Indian baseline is left behind.  The
Sanskrit word \textit{ketu} originally meant ``banner,'' ``bright appearance,''
or ``radiant body,'' and in earlier Vedic and post-Vedic literature it referred
primarily to comets \cite[p.~47]{kotyk2018sinicization}.  The identification of
\textit{ketu} specifically with the descending lunar node is a later development,
consolidated in the \textit{siddh\=anta} tradition but never entirely displacing
the older, broader sense.  This semantic flexibility---a name that had already
migrated once, from comet to node---prefigures the second migration that would
occur in Tang dynasty China, when \zh{計都} moved from node to apogee.

\begin{table*}[ht]
\centering
\small
\caption{Rahu and Ketu: Indian original vs.\ mature Chinese identification.}
\label{tab:rahu-ketu}
\begin{tabular}{@{}lll@{}}
\toprule
Body & Indian (\textit{Jyotish}) & Chinese (\zh{果老星宗}) \\
\midrule
\zh{羅睺} (Rahu) & Ascending node & Ascending node \\
\zh{計都} (Ketu) & Descending node & Osculating lunar apogee \\
\zh{月孛} (Yuebei) & --- & Mean lunar apogee \\
\zh{紫氣} (Ziqi) & --- & Fictitious ($\sim$28-year cycle) \\
\bottomrule
\end{tabular}
\end{table*}

\subsection*{The Clarity That Makes Transformation Visible}

The Indian system's identification of Rahu and Ketu is admirably
unambiguous. Two nodes, always opposite, always moving at the same rate,
always causally connected to eclipses: the conceptual architecture is
self-consistent and empirically grounded.  Ketu is the descending node,
full stop.  There is no competing function for it to perform, no
ambiguity about its period or direction of motion.

Table~\ref{tab:rahu-ketu} places this Indian baseline alongside the
mature Chinese identification that will be traced through the following
sections.  The contrast is arresting.  In the Chinese \zh{七政四餘}
system, \zh{羅睺} (Rahu) retains its Indian identity as the ascending
node, but \zh{計都} has been detached from the descending node
altogether and reattached to the osculating lunar apogee---a point that
moves in the \emph{same} direction as the Moon rather than opposite to
it, has a period of roughly 8.85 years rather than 18.6 years, and has
no direct causal relationship to eclipses.  A new body, \zh{月孛}
(\textit{Yueb\`ei}, ``lunar comet''), has been invented to track the
mean apogee, and a fourth fictitious point, \zh{紫氣} (\textit{Ziq\`{\i}},
``purple vapour''), has been added with a $\sim$28-year cycle that has
no Indian precedent.

This proliferation did not happen overnight.  It was the cumulative
product of successive transformations---transliteration, reidentification,
indigenous restructuring, and hybridization---that will be traced
dynasty by dynasty, beginning with the first moment of contact in
718~CE.

% === §3 The Tang Transmission =================================================
\section{The Tang Transmission: 718--806~CE}\label{sec:tang}
% ~1,500 words — Task 7

\subsection*{Transformation Mode~1: Transliteration}

The first mode of transformation is the simplest: foreign astronomical
concepts are adopted with minimal modification.  Names are transliterated,
tables are copied, procedures are translated as literally as the target
language permits.  The receiving culture does not yet possess the
institutional confidence to reshape what it has received; it is still
in the posture of a student before a prestigious foreign tradition.
This is the characteristic mode of the early Tang reception of Indian
astronomy.

\subsection*{The \textit{Jiuzhi li} of 718~CE}

The pivotal document is the \textit{Jiuzhi li} \zh{九執曆}
(\textit{Navagraha-karana}, ``Calendar of the Nine Seizers''), translated
in 718~CE by Gautama Siddh\=artha (\zh{瞿曇悉達}), an astronomer of
Indian descent whose family had settled in Chang'an over the preceding
generation.  Gautama served in the Tang Bureau of Astronomy
(\zh{太史局}) and brought to his translation both fluency in Sanskrit
astronomical texts and direct access to the imperial court's demand for
accurate calendrical and astrological computation
\cite{yabuuchi1989sui,kotyk2018sinicization}.

The \textit{Jiuzhi li} introduced into China a package of Indian
astronomical innovations that had no precedent in the native tradition.
Most consequential for positional astronomy was the adoption of a
360-degree circle, replacing the Chinese convention of dividing the
ecliptic into 365.25 degrees (one degree per day of the solar year).
Alongside this came decimal notation for fractional positions, sine
tables computed at 3.75° intervals (a characteristic feature of the
Indian \textit{siddh\=anta} tradition that derives ultimately from
Hellenistic chord tables), and epicyclic planetary models for computing
the true positions of the five visible planets from their mean motions.

Among the nine \textit{graha} that give the text its name, two were
purely mathematical points: Rahu (\zh{羅睺}) and Ketu (\zh{計都}).
Their Chinese names are phonetic transliterations---\zh{羅睺} for
``Rahu'' and \zh{計都} for ``Ketu''---the clearest possible signal that
the Tang astronomers were importing the concept as a foreign package
rather than mapping it onto any pre-existing Chinese category.  Both
entered Chinese astronomical vocabulary simultaneously, both designated
lunar nodes, and both were assigned the same westward precessional
motion of approximately 3.06° per month
\cite{yabuuchi1989sui,kotyk2018sinicization}.

The \textit{Jiuzhi li} also introduced the ecliptic as an explicit
reference plane.  Chinese astronomers had for centuries measured
celestial positions in equatorial coordinates: right ascension expressed
as entry into one of the twenty-eight lunar mansions (\zh{入宿度},
\textit{r\`uxiud\`u}) and polar distance (\zh{去極度}, \textit{q\`uj\'ud\`u}).
The ecliptic was not ignored in Chinese astronomy---the twelve
\textit{ci} (\zh{次}) had long served as a solar-position scheme---but
it was not the primary framework for planetary positions.  The
\textit{Jiuzhi li} brought with it a fully ecliptic-centric positional
system, creating an immediate tension between the imported framework and
the native one that would not be resolved for nearly six centuries.

\subsection*{The Iranian Intermediary}

The Tang transmission was not, however, a clean direct line from India
to China.  Kotyk's detailed philological analysis demonstrates that the
\textit{Jiuzhi li} and the astrological texts that accompanied it show
unmistakable signs of Persian and Sasanian intermediary influence
\cite{kotyk2018sinicization,kotyk2017buddhist}.  Calculation procedures,
parameter values, and terminology in several passages align more closely
with Persian \textit{z\={\i}j} tradition than with the Indian
\textit{siddh\=anta} sources from which the system ultimately derives.

Mak's recent study of the Persian astronomical channel into Tang China
strengthens this picture considerably \cite{mak2024persian}.  The
Sasanian empire maintained sophisticated astronomical institutions, and
Persian court astronomers had been assimilating and adapting Indian
methods for at least two centuries before the Tang transmission.  When
Indian astronomical ideas arrived in Chang'an, they had already been
filtered through a Persian cultural lens that introduced modifications---
in nomenclature, in parameter selection, in the relative prominence
accorded different celestial bodies---that are not always easy to
disentangle from the original Indian substrate.

This Iranian intermediary matters for the history of Ketu's
transformation.  If modifications to the node--apogee identification
were already under way in Persian astrology, the Chinese reidentification
may reflect not a purely Chinese innovation but the endpoint of a
gradual drift that began somewhere along the Silk Road between India
and Chang'an.  The evidence is not yet conclusive, but the Persian
channel must remain in view as a possible locus of the first
reidentification.

\subsection*{The \zh{七曜攘災訣} and the Apogee Identification (806~CE)}

The second key Tang text is the \textit{Qiy\`ao r\`angz\=ai ju\'e}
\zh{七曜攘災訣} (``Formulae for Averting Disasters Caused by the Seven
Luminaries''), a Buddhist astrological manual compiled by a
Br\=ahma\d{n}a priest from western India sometime after 806~CE, the
epoch year (Year~1 of the Yuanhe \zh{元和} reign period) embedded in its
planetary tables.  The text would not survive in Chinese repositories;
it was brought to Japan in 865~CE by the Japanese monk Shuy\=u
(\zh{宗叡}) and preserved within the \textit{Sukuy\=od\=o}
(\zh{宿曜道}, ``Way of the Lunar Mansions and Stars'') tradition, where
it remained effectively unknown to Chinese and Western scholarship until
modern times \cite{kotyk2017buddhist}.

The \textit{Qiy\=ao r\`angz\=ai ju\'e} contains multi-decade ephemerides for
all nine \textit{graha}, including a Rahu table spanning 93~years and a
Ketu table spanning 62~years.  In 1994, the Chinese historian of science
Niu Weixing subjected the Ketu ephemeris to a systematic computational
reanalysis, comparing its tabulated positions against two candidates:
the descending lunar node as defined in Indian \textit{siddh\=anta}
astronomy, and the osculating lunar apogee \cite{niu1995ketu}.  The
result was unambiguous: the Ketu ephemeris fits the motion of the
\emph{apogee}, not the node.  The period, the direction of motion, and
the epoch position all correspond to the lunar apogee; none correspond
to the descending node.

This finding establishes a terminus ante quem for the reidentification.
By 806~CE---within roughly one century of the \textit{Jiuzhi li}
translation---\zh{計都} had already ceased to designate the descending
lunar node and had been reassigned to track the osculating apogee.  The
name persisted; the referent changed completely.

\subsection*{Coordinate Thread: The Ecliptic in an Equatorial World}

The coordinate dimension of the Tang transmission is a story of
institutional coexistence rather than replacement.  The \textit{Jiuzhi li}
introduced an ecliptic framework, but the Bureau of Astronomy continued
to produce the official calendar (\textit{li} \zh{曆}) in the native
equatorial system.  Indian ecliptic methods were used for astrological
purposes---planetary hour lords, eclipse prediction, the \textit{graha}
positions in horoscopes---while the equatorial system retained its
authority for the state rituals tied to the solar terms and the
twenty-eight lunar mansions.

The result was a bifurcated institutional landscape.  Buddhist
astrologers working from Indian and Persian sources operated in ecliptic
coordinates; the official Bureau computed in equatorial ones.  The two
systems were not formally reconciled, and conversion procedures between
them were approximate at best.  This bifurcation, established in the
early eighth century, would persist as an unresolved tension through the
Song dynasty, when the first serious attempt at a unified ecliptic
sidereal system for Chinese \textit{qizheng siyu} (\zh{七政四餘},
Seven Governors and Four Remainders) astrology would finally be made.

\subsection*{Identity Thread: Mode~2 Already in Motion}

The \textit{Qiy\=ao r\`angz\=ai ju\'e}'s apogee ephemeris means that by
the end of the Tang period, both transformation modes are already
visible.  Mode~1 (transliteration) is represented by the \textit{Jiuzhi
li}: foreign names and methods are adopted wholesale, with \zh{羅睺}
and \zh{計都} entering Chinese as phonetic loan words for Indian
concepts.  Mode~2 (reidentification) is already under way: the name
\zh{計都} has been retained but silently detached from its original
referent.

What is striking about this early reidentification is its lack of
documentary justification.  No Tang text known to survive explains
\emph{why} Ketu should track the apogee rather than the node.  The
change is not announced, not argued for, not defended.  It appears in
the ephemeris tables as a fait accompli, suggesting that it may have
occurred upstream---in the Persian intermediary tradition, or in a
lost stratum of Indian regional astronomy---rather than being a
conscious Chinese innovation.  Whatever its origin, the reidentification
was consequential: it would shape every subsequent Chinese treatment of
the Four Remainders, creating the distinctive asymmetry---Rahu as node,
Ketu as apogee---that defines the mature \zh{七政四餘} system and
persists in software to this day.

% === §4 The Song System =======================================================
\section{Ecliptic Sidereal: The Song System (\zh{三辰通載})}
\label{sec:song}
% ~1,200 words — Task 8

\textit{Section placeholder.}

% === §5 The Yuan Reform =======================================================
\section{Equatorial Sidereal: The Yuan Reform (\zh{鄭氏星案})}
\label{sec:yuan}
% ~1,500 words — Task 8

\textit{Section placeholder.}

% === §6 The Ming Synthesis ====================================================
\section{Quasi-Ecliptic Sidereal: The Ming Synthesis (\zh{果老星宗})}
\label{sec:ming}
% ~1,800 words — Task 9

\textit{Section placeholder.}

% === §7 The Jesuit Disruption =================================================
\section{Ecliptic Tropical: The Jesuit Disruption (1645)}
\label{sec:jesuit}
% ~2,000 words — Task 9

\textit{Section placeholder.}

% === §8 Algorithmic Fossilization =============================================
\section{Algorithmic Fossilization}\label{sec:fossil}
% ~1,200 words — Task 10

\textit{Section placeholder.}

% === §9 Computational Verification ============================================
\section{Computational Verification}\label{sec:computation}
% ~1,500 words — Task 10

\textit{Section placeholder.}

% === §10 Discussion ===========================================================
\section{Discussion}\label{sec:discussion}
% ~1,200 words — Task 11

\textit{Section placeholder.}

% === §11 Conclusion ===========================================================
\section{Conclusion}\label{sec:conclusion}
% ~500 words — Task 11

\textit{Section placeholder.}

% === Bibliography =============================================================
\bibliographystyle{unsrtnat}
\bibliography{silk-road-astronomy}

\end{document}
